\chapter{Python 環境を整える}

\section{何を分けて考えるべきか}

Python を使い始めるときは、次の 3 つを混同しない方がよい。

\begin{enumerate}
\item Python 本体を入れること
\item 複数版の Python を切り替えること
\item プロジェクトごとの仮想環境を作ること
\end{enumerate}

まず Python 本体が必要であり、そのうえで Python 3.11 と 3.12 を切り替えたいなら版管理ツールが必要になる。さらに、同じ Python 3.12 でもプロジェクト A と B で依存ライブラリが違うなら、仮想環境で分ける必要がある。この 3 段階を分けて理解した方が、あとで混乱しにくい。

\section{今の Python を確認する}

macOS でも WSL でも、最初にいま見えている Python を調べておく方がよい。

\begin{lstlisting}[numbers=none, xleftmargin=2\zw, caption={Python の確認}]
$ python3 --version
$ which python3
$ python3 -c "import sys; print(sys.executable)"
\end{lstlisting}

\code{python3 --version} で版を確認し、\code{which python3} と \code{sys.executable} で実体の場所を確認する。複数版を切り替えるようになると、この確認は非常に重要になる。

\section{\code{pyenv} で複数版の Python を管理する}

Python 本体の導入と切り替えには \code{pyenv} が分かりやすい。\code{pyenv} は、シェルから見える \code{python} をディレクトリ単位で切り替えるための道具であり、複数版の Python を使い分けるときの基準として扱いやすい。

\subsection{macOS での準備}

\code{pyenv} は Python 本体をビルドして入れるので、macOS では Xcode Command Line Tools が必要になる。

\begin{lstlisting}[numbers=none, xleftmargin=2\zw, caption={Command Line Tools があるか確認する}]
$ xcode-select -p
\end{lstlisting}

すでに Homebrew が正常に動いているなら、Command Line Tools は入っていることが多い。上のコマンドでパスが表示されれば、そのまま進んでよい。表示されない場合は次を実行する。

\begin{lstlisting}[numbers=none, xleftmargin=2\zw, caption={Command Line Tools を入れる}]
$ xcode-select --install
\end{lstlisting}

\subsection{WSL / Ubuntu での準備}

WSL の Ubuntu では、ビルドに必要なライブラリを先に入れておく。

\begin{lstlisting}[numbers=none, xleftmargin=2\zw, caption={Ubuntu での準備}]
$ sudo apt update
$ sudo apt install -y build-essential curl git libssl-dev zlib1g-dev \
  libbz2-dev libreadline-dev libsqlite3-dev libffi-dev liblzma-dev tk-dev
\end{lstlisting}

\subsection{\code{pyenv} を入れる}

\begin{lstlisting}[numbers=none, xleftmargin=2\zw, caption={pyenv の導入}]
$ curl -fsSL https://pyenv.run | bash
\end{lstlisting}

\subsection{どの設定ファイルを編集するか決める}

\code{pyenv} を使うには、シェルの設定ファイルへ初期化コードを書く必要がある。どのファイルを編集するかは、いま使っているシェルで決まる。

\begin{lstlisting}[numbers=none, xleftmargin=2\zw, caption={シェルを調べる}]
$ echo $SHELL
\end{lstlisting}

\code{/bin/zsh} と出たら \code{\textasciitilde/.zshrc}、\code{/bin/bash} と出たら \code{\textasciitilde/.bashrc} を編集する。これで、\code{zsh} や \code{bash} が唐突に出てくる状態は避けられる。

\subsection{\code{pyenv} をシェルへ組み込む}

\begin{lstlisting}[numbers=none, xleftmargin=2\zw, caption={zsh での設定}]
$ echo 'export PYENV_ROOT="$HOME/.pyenv"' >> ~/.zshrc
$ echo '[[ -d $PYENV_ROOT/bin ]] && export PATH="$PYENV_ROOT/bin:$PATH"' >> ~/.zshrc
$ echo 'eval "$(pyenv init - zsh)"' >> ~/.zshrc
$ exec "$SHELL"
\end{lstlisting}

\begin{lstlisting}[numbers=none, xleftmargin=2\zw, caption={bash での設定}]
$ echo 'export PYENV_ROOT="$HOME/.pyenv"' >> ~/.bashrc
$ echo '[[ -d $PYENV_ROOT/bin ]] && export PATH="$PYENV_ROOT/bin:$PATH"' >> ~/.bashrc
$ echo 'eval "$(pyenv init - bash)"' >> ~/.bashrc
$ exec "$SHELL"
\end{lstlisting}

\subsection{版を入れて切り替える}

\code{pyenv} を使う場合は、\code{pyenv install} で版を追加し、\code{pyenv global}、\code{pyenv local}、\code{pyenv shell} で切り替える。

\begin{lstlisting}[numbers=none, xleftmargin=2\zw, caption={版を追加して確認する}]
$ pyenv install 3.11.11
$ pyenv install 3.12.10
$ pyenv versions
\end{lstlisting}

\begin{lstlisting}[numbers=none, xleftmargin=2\zw, caption={普段使う版を決める}]
$ pyenv global 3.12.10
$ python --version
$ which python
\end{lstlisting}

\code{pyenv global} は普段使う版を決める。これに対して \code{pyenv local} は、そのディレクトリだけ別の版を使いたいときに使う。

\begin{lstlisting}[numbers=none, xleftmargin=2\zw, caption={プロジェクト単位で版を固定する}]
$ mkdir analysis311
$ cd analysis311
$ pyenv local 3.11.11
$ python --version
$ cat .python-version
\end{lstlisting}

\code{pyenv local} を実行すると \code{.python-version} が作られ、そのディレクトリでは指定した版が優先される。これによって、同じマシンの中でプロジェクト A は 3.11、プロジェクト B は 3.12 という使い分けができる。

\code{pyenv shell 3.11.11} は、そのターミナルを閉じるまでの一時的な切り替えである。検証用に 1 回だけ古い版で動かしたいときに便利である。

\section{\code{venv} でプロジェクトごとの仮想環境を作る}

Python 本体の版を決めたら、その版の上にプロジェクト専用の仮想環境を作る。依存ライブラリの衝突を避けるためである。

\begin{lstlisting}[numbers=none, xleftmargin=2\zw, caption={venv の基本}]
$ python -m venv .venv
$ source .venv/bin/activate
$ which python
$ which pip
$ python --version
$ pip --version
\end{lstlisting}

仮想環境へ入ると、プロンプトの先頭に \code{(.venv)} が付くことが多い。

\begin{lstlisting}[numbers=none, xleftmargin=2\zw, caption={プロンプトの例}]
(.venv) ikomae@macbook analysis $
\end{lstlisting}

\code{which python} と \code{which pip} を見ると、\code{.venv/bin/python} と \code{.venv/bin/pip} が使われていることが分かる。これは、\code{source .venv/bin/activate} が \code{\$PATH} の先頭に \code{.venv/bin} を追加しているからである。

\subsection{\code{pip} と \code{python -m pip}}

仮想環境へ入っているなら、通常は \code{pip} と \code{python -m pip} は同じ環境の \code{pip} を指す。したがって、作業中は \code{pip install numpy} と書いてよい。

一方で、いまどの Python が有効か不安なときや、スクリプトや文書の中で「この Python に対応する \code{pip} を使う」と明示したいときは、\code{python -m pip} の方が誤解が少ない。実務では両方見るが、意味の違いを分かって使うことが重要である。

\begin{lstlisting}[numbers=none, xleftmargin=2\zw, caption={仮想環境での導入例}]
$ pip install --upgrade pip
$ pip install numpy pandas matplotlib scipy astropy tqdm
\end{lstlisting}

作業を終えたら \code{deactivate} で仮想環境から抜けられる。

\section{JupyterLab を併用する}

対話的に試しながら解析したいときは、JupyterLab も便利である。小さなコード片をその場で動かし、数値や図をすぐ確認できるため、探索的な作業では特に役立つ。

\begin{lstlisting}[numbers=none, xleftmargin=2\zw, caption={JupyterLab の導入と起動}]
$ pip install jupyterlab
$ jupyter lab
\end{lstlisting}

JupyterLab の利点は、コードを実行できるだけでなく、その前後に説明文、数式、図の読み方を書き添えられる点にある。解析の途中経過を他人へ見せるときも、自分で後から読み返すときも、コードと説明が同じノートブックに並んでいる方が理解しやすい。

\begin{lstlisting}[numbers=none, xleftmargin=2\zw, caption={Markdown セルの例}]
# 解析メモ

- 入力ファイル: `events_2024_01_01.txt.gz`
- 閾値: `signal > 5.0`

$\Delta t$ の分布を確認する
\end{lstlisting}

ただし、JupyterLab は最終成果物の置き場所ではなく、試行錯誤の場として使う方が整理しやすい。解析手順を再現したい段階では、ノートブック中で固めた処理を Python スクリプトへ戻し、入力や出力を明示した方が後で困りにくい。

\section{\code{uv} はプロジェクト管理に使う}

\code{uv} は Python 本体の導入もできるが、本書ではそれよりも「プロジェクト単位の管理」に使う道具として位置付ける方が自然である。特に、依存関係、ロックファイル、\code{.python-version}、Git 管理、パッケージ化、ワークスペース管理といった、プロジェクト全体の整理で力を発揮する。

\subsection{まずは既定の \code{uv init} を使う}

多くの人は、まずシンプルに \code{uv init} を使えばよい。オプションは後から必要になったときに足せばよい。

\begin{lstlisting}[numbers=none, xleftmargin=2\zw, caption={uv で新しいプロジェクトを始める}]
$ mkdir analysis_uv
$ cd analysis_uv
$ uv init
\end{lstlisting}

既定の \code{uv init} では、少なくとも次のようなファイルやディレクトリが作られる。

\begin{lstlisting}[numbers=none, xleftmargin=2\zw, caption={既定の uv init が作るもの}]
.git/
.gitignore
.python-version
README.md
main.py
pyproject.toml
\end{lstlisting}

つまり、\code{uv init} は単に Python ファイルを 1 つ置くだけではなく、Git 管理の土台までまとめて用意する。この点が \code{venv} 単体との大きな違いである。

\subsection{必要ならオプションを使う}

既定動作を変えたいときだけオプションを使えばよい。

\begin{itemize}
\item \code{uv init --bare}: \code{pyproject.toml} だけを作る
\item \code{uv init --vcs none}: Git 初期化をしない
\item \code{uv init --no-pin-python}: \code{.python-version} を作らない
\item \code{uv init --package} や \code{uv init --lib}: スクリプトではなくパッケージやライブラリとして始める
\end{itemize}

特に、コードを再利用可能なモジュールやライブラリとして育てたいなら、\code{--package} や \code{--lib} を意識した方がよい。単発スクリプトで終わるのか、将来 \code{import} される部品になるのかで、適切な初期構成は変わる。

\subsection{依存関係を追加する}

\code{uv add} は依存ライブラリを追加し、\code{pyproject.toml} とロック情報を更新する。

\begin{lstlisting}[numbers=none, xleftmargin=2\zw, caption={依存関係を追加する}]
$ uv add numpy pandas matplotlib scipy astropy tqdm
\end{lstlisting}

\subsection{環境をそろえる}

\code{uv sync} は、\code{pyproject.toml} とロックファイルに基づいて環境をそろえる。ここで初めて \code{.venv} が作られることもある。

\begin{lstlisting}[numbers=none, xleftmargin=2\zw, caption={環境を同期する}]
$ uv sync
\end{lstlisting}

過去の自分の解析や、他人から受け取ったプロジェクトを再現するときは、この \code{uv sync} が重要になる。リポジトリを取得して \code{uv sync} すれば、必要な依存関係をまとめてそろえやすい。ここで Git の基本が必要になるので、付録~\ref{app:git} もあわせて読むとよい。

\subsection{コマンドを実行する}

\code{uv run} は、そのプロジェクトの環境でコマンドを実行する。スクリプトを直接走らせるときと、Python 本体を起動したいときとで書き分ける。

\begin{lstlisting}[numbers=none, xleftmargin=2\zw, caption={uv run の例}]
$ uv run main.py
$ uv run python --version
$ uv run python -m pytest
\end{lstlisting}

実行対象が Python スクリプトそのものであるなら、\code{uv run main.py} と書けばよい。わざわざ \code{python} を挟む必要はない。一方で、\code{-m} を使いたいときや、Python 本体へオプションを渡したいときは \code{uv run python ...} と書く。

\subsection{Python 版の固定にも使える}

\code{uv} もプロジェクト単位で Python 版を固定できる。\code{pyenv} がシェル全体の切り替えを担当し、\code{uv} がプロジェクト側の固定を担当する、と分けて考えると整理しやすい。

\begin{lstlisting}[numbers=none, xleftmargin=2\zw, caption={uv で Python 版を固定する}]
$ uv python pin 3.12
$ uv run python --version
$ uv run --python 3.11 python --version
\end{lstlisting}

\section{標準ライブラリと外部ライブラリ}

Python には、最初から使える標準ライブラリと、自分で追加する外部ライブラリがある。

\begin{itemize}
\item 標準ライブラリの例: \code{math}, \code{pathlib}, \code{glob}, \code{argparse}
\item 外部ライブラリの例: \code{numpy}, \code{pandas}, \code{matplotlib}, \code{scipy}, \code{astropy}
\end{itemize}

標準ライブラリは Python 本体に含まれている。外部ライブラリは、目的に応じて自分で追加する。解析では外部ライブラリが非常に重要だが、まずは「何が標準で、何が追加なのか」を意識しておくと整理しやすい。

\section{解析用のライブラリを入れる}

本書の後半で使う主なライブラリを表~\ref{tab:packages} に示す。

\begin{table}[hbtp]
\centering
\caption{本書で使う主なライブラリと役割}
\label{tab:packages}
\begin{tabular}{ll}
\toprule
ライブラリ & 主な役割 \\
\midrule
NumPy & 数値配列の高速処理 \\
pandas & 表形式データの読込みと加工 \\
Matplotlib & 図の作成 \\
SciPy & フィットと数値計算 \\
Astropy & 時刻・単位・座標の扱い \\
tqdm & 進捗表示 \\
\bottomrule
\end{tabular}
\end{table}

\code{venv} を使うなら \code{pip install}、\code{uv} を使うなら \code{uv add} と書けばよい。

\begin{lstlisting}[numbers=none, xleftmargin=2\zw, caption={venv を使う場合}]
$ pip install numpy pandas matplotlib scipy astropy tqdm
\end{lstlisting}

\begin{lstlisting}[numbers=none, xleftmargin=2\zw, caption={uv を使う場合}]
$ uv add numpy pandas matplotlib scipy astropy tqdm
\end{lstlisting}

\section{コマンドライン引数と設定ファイルを使い分ける}

同じ解析スクリプトを何度も使うなら、入力ディレクトリや出力先をコードに埋め込むより、コマンドライン引数で渡す方がよい。Python では \code{argparse} が標準で使える。

\begin{lstlisting}[language=Python, caption={argparse の最小例}]
import argparse

parser = argparse.ArgumentParser()
parser.add_argument("input_dir")
parser.add_argument("--outdir", default="output")
args = parser.parse_args()

print(args.input_dir, args.outdir)
\end{lstlisting}

ただし、何でも引数にすればよいわけではない。毎回ほぼ同じ値まで毎回書かせると、使いにくいスクリプトになる。

値の置き場所は、表で機械的に覚えるより、用途ごとに言葉で整理した方が分かりやすい。

\begin{description}
\item[コマンドライン引数] 実行ごとに変わる値を置く。入力ディレクトリ、出力先、日付範囲、並列数などが典型である。
\item[設定ファイル] プロジェクト単位では固定だが、将来変更し得る値を置く。較正ファイルのパス、検出器 ID、閾値、図の保存形式などが当てはまる。
\item[環境変数] マシン依存の値や秘密情報を置く。\code{DATA\_ROOT}、API トークン、認証情報などが典型である。
\item[コード内の定数] 定義として固定したい値や単位変換を書く。\code{NS\_PER\_SECOND} や物理定数のように、意味が仕様として決まっているものを置く。
\end{description}

入力ファイルのパス、出力先、しきい値、作業ディレクトリのように「実行条件」であるものは、原則としてハードコーディングしない方がよい。コードを書き換えずに切り替えられるようにしておくと、再利用と再現が楽になる。

\subsection{設定ファイルを読む}

めったに変わらない値は、設定ファイルへまとめる方が使いやすい。形式は JSON でも TOML でもよい。ここでは標準ライブラリだけで扱える JSON の例を示す。

\begin{lstlisting}[numbers=none, xleftmargin=2\zw, caption={設定ファイルの例}]
{
  "calibration_path": "calib/detector_a.csv",
  "coincidence_window_ns": 250,
  "figure_format": "pdf"
}
\end{lstlisting}

\begin{lstlisting}[language=Python, caption={JSON 設定を読み込む例}]
import json
from pathlib import Path

config = json.loads(Path("config.json").read_text())
window_ns = config["coincidence_window_ns"]
\end{lstlisting}

このようにしておけば、同じコードを保ったまま設定だけを差し替えられる。Python 3.11 以降なら \code{tomllib} を使って TOML を読むこともできる。どの形式を選ぶにしても、「たまに変わるがコード本体ではない値」を分離することが大切である。

\subsection{マジックナンバーを避ける}

\code{250} や \code{3600} のような数値が説明なしにコードへ裸で出てくると、後から読んだ人は意味を判断しにくい。このような数値はマジックナンバーと呼ばれる。詳しくは \ref{sec:magic-number} 節で述べるが、意味のある数値には名前を付け、実行条件なら引数や設定ファイルへ出す、という整理を早い段階で身に付けた方がよい。
